\documentclass[11pt]{article}
\usepackage[review]{acl}   % 'review' = numeración de líneas + anonimato. Camera-ready: \usepackage[]{acl}
\usepackage{times}
\usepackage{latexsym}
\usepackage[T1]{fontenc}
\usepackage[utf8]{inputenc}
\usepackage{microtype}
\usepackage{graphicx}
\usepackage{booktabs}
\usepackage{amsmath}
\usepackage{amssymb}

% ===========================================================================================
% NOMBRE DEL MODELO: TRACE (TRajectory-Aware Contextual Encoder) -- NOMBRE DE TRABAJO, NO DEFINITIVO.
%   - sigla limpia: TRajectory-Aware Contextual Encoder. Anclado en la línea: D2F lee los diálogos
%     como "trayectorias en un espacio latente"; el paper de ANN es "Retrieving the Flow".
%     TRACE = la traza/trayectoria de la conversación. Título en SERIE: "Anticipating the Flow".
%   - Considerado y descartado por ahora: TRAZAR (esp. "trazar", -AR x Argentina/autoregresivo) —
%     no convenció el acrónimo. Alternativas vivas: HORIZON, CURRENT, FLUX.
%   - TODO(JM): el nombre se revisita; chequear colisión (TRACE = modelo viejo de psicolingüística,
%     McClelland & Elman 1986, otro dominio). Cambiarlo = find-replace, no bloquea nada.
% ===========================================================================================

\title{Anticipating the Flow: Learned Contextual Turn Embeddings\\ for Dialogue-Act Trajectory in Task-Oriented Dialogue}

% TODO(JM): lista/orden de autores a confirmar con los directores. En 'review' va anónimo igual.
\author{Juan Manuel Fern\'andez \\
  Universidad Nacional de Luj\'an / UNSL \\
  \texttt{jmfernandez@unlu.edu.ar} \\\And
  Sergio Burdisso \\
  Idiap Research Institute \\\And
  Marcelo Errecalde \\
  Universidad Nacional de San Luis}

\begin{document}
\maketitle

\begin{abstract}
In task-oriented dialogue, turns are commonly represented as dense embeddings produced
independently for each utterance. Such per-turn representations capture the local content of a
turn but not its position within the conversational trajectory, which depends on the accumulated
history and on where the dialogue is heading. We introduce \textbf{TRACE} (TRajectory-Aware
Contextual Encoder), a pre-trained contextual turn encoder that, on top of a frozen per-turn
embedding space (Dialog2Flow), learns in a self-supervised manner to integrate the dialogue
context into each turn representation, in both autoregressive and bidirectional formulations. The
construction is analogous to contextual token representations in language models, but operates at
the level of dialogue turns rather than tokens.

We evaluate whether TRACE encodes conversational trajectory through a frozen linear probe, in the
tradition of feature-based evaluation of pre-trained encoders: the representation is held fixed and
a logistic-regression classifier is trained to predict dialogue acts. We distinguish the
\emph{current} act of a turn---available in the per-turn embedding---from the act of the
\emph{next} turn, whose prediction requires the conversational trajectory. On held-out dialogues
that the encoder never saw during pre-training, TRACE improves next-act prediction over the frozen
per-turn base, over hand-crafted cumulative and EMA memories, and over a context-aware dialogue
model (TOD-BERT), by $0.12$--$0.14$ macro-F1.

These results indicate that a learned contextual turn encoder captures act-trajectory information
that per-turn encoders and fixed accumulation policies do not.
\end{abstract}

\section{Introduction}
\label{sec:intro}

Task-oriented dialogue (TOD) systems interact with a user to complete a goal, such as booking a
restaurant or querying an external service. To do so, the agent must interpret each utterance,
maintain a coherent representation of the dialogue state, and select an appropriate action at every
turn~\cite{young2013pomdp,chen2017survey}. A common way to represent dialogue turns is through
dense embeddings produced by pre-trained encoders, which enable retrieval, classification, and
clustering over large conversational collections~\cite{devlin2019bert,reimers2019sentencebert}.

Dialogues are, however, inherently sequential: the meaning of a turn depends on the accumulated
context and on its role in the conversational flow. Encoding each turn in isolation may therefore
fail to capture how the dialogue state evolves, especially when the relevant information is spread
across several turns. This observation has motivated representations that integrate the history,
either through fixed accumulation policies or through latent models of the dialogue
state~\cite{wen2018sequence,min2020dialogue,hudecek2022latent,liu2023vls}.

In a companion study~\cite{fernandez2025retrieving} we examined post-hoc accumulation policies
---static, normalized cumulative, and exponential moving average (EMA) representations--- over
several embedding spaces, and found that their usefulness for conversational-memory retrieval
depends on the geometry of the base space and on the update policy. That work closed by pointing to
a natural next step: replacing the \emph{fixed} update rule with a \emph{learned} one, so that each
turn representation is produced as a function not only of the turn itself but also of its
conversational context. \textbf{The present work takes that step.}

We introduce \textbf{TRACE} (TRajectory-Aware Contextual Encoder), a pre-trained contextual turn
encoder. TRACE operates on top of a frozen per-turn embedding space---Dialog2Flow
(D2F)~\cite{burdisso2024dialog2flow}, a representation trained to organize utterances by their
communicative function---and learns, through self-supervised objectives over entire dialogues, an
update function that contextualizes each turn. The design mirrors contextual token representations
in language models~\cite{devlin2019bert}, but at the granularity of dialogue turns; we consider both
an autoregressive variant, conditioned only on preceding turns, and a bidirectional one, conditioned
on the whole dialogue. We frame TRACE as an extension of the Dialog2Flow line, not as a competitor:
D2F provides the per-turn substrate, and TRACE adds the contextual layer above it.

Our central question is whether such a learned encoder captures conversational trajectory that
per-turn encoders lack. We address it with a \emph{frozen linear probe}, in the tradition of
feature-based evaluation of pre-trained representations~\cite{devlin2019bert}: we hold the
representation fixed and train only a linear classifier to predict dialogue acts. We separate the
\emph{current} act of a turn---a control task, since this information already lives in the per-turn
embedding---from the act of the \emph{next} turn, whose prediction genuinely requires the dialogue
trajectory. On held-out dialogues never seen during pre-training, TRACE improves next-act prediction
over the frozen per-turn base, over hand-crafted cumulative and EMA memories, and over a
context-aware dialogue model, indicating that the contextual layer encodes trajectory information
that the alternatives do not.

The remainder of the paper is organized as follows. Section~\ref{sec:background} reviews related
work. Section~\ref{sec:method} describes the TRACE architecture and its pre-training objectives.
Section~\ref{sec:setup} presents the evaluation protocol, and Section~\ref{sec:results} the results.
Section~\ref{sec:discussion} discusses their scope and limitations, and Section~\ref{sec:conclusion}
concludes.

\section{Background and Related Work}
\label{sec:background}

\paragraph{Turn and sentence embeddings.}
Pre-trained encoders map utterances to dense vectors that capture semantic content and support
similarity, search, and clustering~\cite{reimers2019sentencebert}. General-purpose sentence encoders
such as those based on MiniLM~\cite{wang2020minilm} and MPNet~\cite{song2020mpnet} are widely used,
but they encode each turn independently and are not specific to dialogue. Applied per turn, they
capture local content rather than the position of the turn within a broader conversational
trajectory.

\paragraph{Contextual representations.}
Transformer-based language models produce \emph{contextual} token
representations~\cite{vaswani2017attention,devlin2019bert}: the representation of a token depends on
the surrounding sequence. TRACE transposes this idea to dialogue, treating turns as the sequence
elements and producing a representation of each turn that depends on the dialogue. Unlike token-level
pre-training, however, the input units here are continuous per-turn embeddings rather than discrete
tokens, which has consequences for the training objective (Section~\ref{sec:method}).

\paragraph{Dialogue-specific representations.}
TOD-BERT~\cite{wu2020todbert} adapts BERT to task-oriented dialogue by pre-training over concatenated
dialogue context with speaker tokens; it is context-aware, but it re-encodes raw text rather than
building on a per-turn space. Dialog2Flow~\cite{burdisso2024dialog2flow} pre-trains
action-driven sentence embeddings that organize utterances by communicative function and read
dialogues as trajectories in a latent space, unifying multiple corpora---including
MultiWOZ and Schema-Guided Dialogue---under a common annotation. We take D2F as the frozen per-turn
substrate for TRACE. A parallel line induces latent dialogue states or actions with limited
supervision~\cite{wen2018sequence,min2020dialogue,hudecek2022latent,liu2023vls}; TRACE differs in
that it learns a contextual \emph{representation} of turns through self-supervision, rather than an
explicit latent-state model.

\paragraph{Probing pre-trained encoders.}
A standard way to assess what a frozen representation encodes is to train a simple classifier on top
of it for a downstream task, keeping the representation
fixed~\cite{devlin2019bert}. We adopt this feature-based protocol, which isolates the quality of the
representation from the capacity of the classifier, and apply it to dialogue-act prediction.

\section{TRACE}
\label{sec:method}

\subsection{Two-stage architecture}
TRACE follows a two-stage pipeline. A frozen per-turn encoder $f_1$ maps each utterance $u_t$ to a
base embedding,
\begin{equation}
  e_t = f_1(u_t), \qquad e_t \in \mathbb{R}^d,
\end{equation}
and a learned contextual encoder $f_2$ maps the sequence of base embeddings of a dialogue to a
sequence of contextual turn representations,
\begin{equation}
  (h_1,\dots,h_T) = f_2(e_1,\dots,e_T), \qquad h_t \in \mathbb{R}^d.
\end{equation}
In this work $f_1$ is \texttt{dialog2flow-joint-bert-base}~\cite{burdisso2024dialog2flow} and is kept
\emph{frozen}: this isolates the contribution of $f_2$ and makes $h_t$ directly comparable to $e_t$.

The contextual encoder $f_2$ is a Transformer~\cite{vaswani2017attention} over the sequence of turn
embeddings of a dialogue. Each $e_t$ is augmented with a learned position embedding and a learned
speaker embedding that encodes the role of the participant, and a stack of self-attention layers lets
every turn attend over the dialogue. We consider two attention regimes: a \emph{bidirectional} one, in
which a turn attends to the entire dialogue, and an \emph{autoregressive} (causal) one, in which it
attends only to preceding turns. Because the sequence elements are continuous per-turn embeddings
rather than discrete tokens, $f_2$ ingests $e_t$ directly---through a linear projection when the base
and hidden dimensions differ---rather than through a token-embedding lookup; this is the main
structural difference from a token-level encoder such as BERT~\cite{devlin2019bert}.

\subsection{An interchangeable per-turn base}
A defining property of the architecture is that $f_1$ is not tied to a particular encoder: $f_2$
operates on the \emph{output space} of $f_1$, so any per-turn encoder can serve as the base. The input
dimension is read from $f_1$, the linear projection above adapts it to the hidden size, and the
residual formulation (below) keeps $h_t$ in the same space as $e_t$; no component of $f_2$ is bound to
a specific base. TRACE is therefore better understood as a \emph{method}---a contextual layer over an
arbitrary per-turn space---than as a single model bound to one encoder. We instantiate it on
Dialog2Flow, which provides a strong action-driven base, and report in Section~\ref{sec:results} that
the contextual gain persists when $f_1$ is replaced by a general-purpose encoder (MPNet) and by a
context-aware dialogue model (TOD-BERT). A systematic study of which base best supports the contextual
layer is left for future work.

\subsection{From fixed to learned accumulation}
The companion study~\cite{fernandez2025retrieving} produced one state vector per turn through a fixed
update, the normalized cumulative representation $h_t = \mathrm{LayerNorm}(h_{t-1} + e_t)$ and its EMA
generalization. TRACE replaces this fixed rule with a learned function of the whole context. To keep
$h_t$ anchored to the base space---so that it remains comparable to $e_t$ and usable in the same
geometry---we use a residual formulation,
\begin{equation}
  h_t = \mathrm{LayerNorm}\!\left(e_t + \Delta_t\right),
  \qquad \Delta_t = f_2^{\Delta}(e_1,\dots,e_T)_t,
\end{equation}
where $\Delta_t$ is the learned contextual update~\cite{ba2016layer}. The fixed cumulative and EMA
policies are thus special cases of the function class that $f_2$ can represent.

\subsection{Self-supervised objectives}
TRACE is pre-trained without functional labels (no dialogue acts, slots, or intents are used during
pre-training). We consider two formulations. In the \emph{bidirectional} setting, each turn is
conditioned on the entire dialogue and the objective reconstructs masked turns. In the
\emph{autoregressive} setting, each turn is conditioned only on preceding turns and the objective
predicts the next turn. Both are complemented by an in-batch contrastive term that scores the correct
target turn against other turns in the batch, treating turn embeddings as a ``vocabulary'' of turns.
% TODO(JM): detalles de pre-entrenamiento (capas/heads/optimizador/épocas/colección) en una
% subsección o apéndice; decidir si presentamos UN modelo (recomendado) o el lineage de variantes.

\section{Experimental Setup}
\label{sec:setup}

\subsection{Data and held-out protocol}
We use the experimental collection of the companion study, derived from
Dialog2Flow~\cite{burdisso2024dialog2flow}: $101{,}021$ complete dialogues and $1{,}000{,}023$ turns.
To measure \emph{inductive} generalization rather than memorization, we hold out a fixed set of
$17{,}362$ dialogues (seed $42$) that are \emph{excluded} from pre-training, and report all probe
results on this held-out split. Untrained baselines are unaffected by this split, so the comparison
isolates generalization.

\subsection{Task: dialogue-act anticipation}
Each turn carries a dialogue-act label. We define two probing tasks. The \emph{current-act} task,
$\mathrm{act}(t)$, predicts the act of the turn itself and serves as a control: this information is
already present in the per-turn embedding, since D2F is action-driven. The \emph{next-act} task,
$\mathrm{act}(t{+}1)$, predicts the act of the following turn and is the actual test of trajectory,
as it cannot be read off the current turn in isolation.

\subsection{Frozen linear probe and baselines}
For each representation we freeze it and train a logistic-regression
classifier~\cite{devlin2019bert}; we report macro-F1 (the act classes are imbalanced) and accuracy.
We compare a ladder of representations, each isolating one factor: a random-feature floor; a
general-purpose sentence encoder (SBERT/MPNet)~\cite{reimers2019sentencebert,song2020mpnet}; the
frozen per-turn base $e_t$ (D2F); hand-crafted memories over $e_t$ (normalized cumulative and EMA);
and an external context-aware model, TOD-BERT~\cite{wu2020todbert}. The autoregressive TRACE variant
is the fair comparison for next-act prediction, since the bidirectional variant has access to the
target turn within its window.
% TODO(JM): significancia (Wilcoxon/bootstrap por diálogo), como en el paper de ANN; breakdown
% por dataset; tamaño del split de probe.

\section{Results}
\label{sec:results}

Table~\ref{tab:main} reports next-act prediction on the held-out dialogues. Representations without
learned turn context---random features, the general-purpose encoder, the per-turn base $e_t$, and the
hand-crafted cumulative and EMA memories---fall in a band up to $0.58$ macro-F1. The external
context-aware model TOD-BERT reaches $0.57$, on par with the per-turn base. TRACE (autoregressive)
attains $0.71$ macro-F1, an improvement of $0.12$--$0.14$ over the strongest per-turn and
hand-crafted baselines and over TOD-BERT, on dialogues that the encoder never saw during
pre-training. Since these baselines are untrained or hand-crafted, and the split is held out, the gap
reflects trajectory information learned by the contextual layer rather than memorization.
% TODO(JM): +significancia (p-values), +fila act(t) de control (debería quedar pareja y alta -> TRACE
% no pierde el acto actual), +sub-sección de robustez agnóstica a la base.

\begin{table}[t]
\centering\small
\begin{tabular}{lcc}
\toprule
Representation & Acc. & Macro-F1 \\
\midrule
Random                         & 0.421 & 0.060 \\
Cumulative                     & 0.577 & 0.458 \\
SBERT (general-purpose)        & 0.623 & 0.511 \\
TOD-BERT (context-aware)       & 0.676 & 0.572 \\
$e_t$ (D2F, frozen base)       & 0.679 & 0.567 \\
EMA ($\alpha{=}0.6$)           & 0.685 & 0.582 \\
\textbf{TRACE (autoregressive)} & \textbf{0.792} & \textbf{0.706} \\
\bottomrule
\end{tabular}
\caption{Next-act prediction ($\mathrm{act}(t{+}1)$) on \textbf{held-out} dialogues (inductive),
frozen linear probe, macro-F1 and accuracy. TRACE outperforms the frozen per-turn base, hand-crafted
memories, and the context-aware TOD-BERT.}
\label{tab:main}
\end{table}

\paragraph{Robustness to the base encoder.}
Because $f_1$ is interchangeable (Section~\ref{sec:method}), we ask whether the contextual gain is
specific to Dialog2Flow. We re-train $f_2$ on top of two qualitatively different bases---a
general-purpose encoder (MPNet) and a context-aware dialogue model (TOD-BERT, used as a single-turn
encoder)---and find that the contextual layer improves next-act prediction over the corresponding base
in each case, indicating that the effect is a property of the method rather than of one base.
% TODO(JM): tabla con D2F/MPNet/TOD-BERT al MISMO protocolo (in-domain held-out) + significancia.
% Hoy el dato disponible es a escala reducida; correrlo al protocolo principal.

\section{Discussion}
\label{sec:discussion}

% NOTA(JM): sacamos el contraste IR/retrieval — era tarea equivocada (TRACE no es para IR) y
% mostrarlo como "no mejora retrieval" era autodebilidad. El retrieval es territorio del paper
% companion ("Retrieving the Flow"); el negativo de los contextuales-para-IR, si se reporta, va ahí.
% TODO(JM): la Discusión puede crecer con la robustez agnóstica a la base.

\paragraph{The ceiling of next-act prediction.}
Anticipating the next act is intrinsically harder than reading the current one: it is a forecasting
problem with irreducible uncertainty (several acts may validly follow a given history) and with label
noise (turns may carry multiple acts; annotations vary across the unified corpora). A bidirectional
variant that can see the target turn within its window reaches only $\approx 0.89$ macro-F1, not
$1.0$, which bounds the attainable performance from above; the reported gains should therefore be
read against the reachable headroom rather than against a perfect score.

\section{Conclusion and Future Work}
\label{sec:conclusion}

We introduced TRACE, a pre-trained contextual turn encoder for task-oriented dialogue that learns,
on top of a frozen per-turn space, to integrate the conversational trajectory into each turn. Using a
frozen linear probe over held-out dialogues, we found that TRACE encodes next-act information that
per-turn encoders, hand-crafted memories, and an external context-aware model do not. This positions
TRACE as the learned counterpart to the fixed accumulation policies of the companion study, and as a
contextual layer that extends the Dialog2Flow line.

Several directions follow. A self-supervised discrete objective---predicting cluster identities from a
codebook derived from the data (e.g., vector quantization or Dialog2Flow flow states), rather than
from human labels---would provide a BERT-like training signal while preserving the ability to scale
to unlabeled text. Joint fine-tuning of $f_1$ and $f_2$, and a systematic study of which per-turn base
encoder best supports the contextual layer, are also natural extensions. Finally, evaluating TRACE on
concrete downstream tasks such as next-action selection would test its usefulness beyond probing.

\section*{Limitations}
% (ACL requiere esta sección.)
The per-turn encoder $f_1$ is kept frozen, so the quality of the contextual representation is bounded
by that of the base space. Our evaluation is in-domain: held-out dialogues are drawn from the same
collection used for pre-training, and cross-corpus and human-human external transfer remain open.
The probe is frozen (a diagnostic of representation quality), not a fine-tuned, competitive system.
% TODO(JM): agregar significancia y, si se hace, transferencia leave-one-dataset-out limpia.

\nocite{*}  % esqueleto: incluye todas las refs aunque algunas \cite estén pendientes. Quitar al cerrar.
\bibliography{references}

\end{document}
